\chapter{Introducción}

\section{Mercado eléctrico}

El Sistema Eléctrico Nacional (SEN) de Chile ha experimentado una transformación estructural sin precedentes durante la última década. La unificación del Sistema Interconectado del Norte Grande (SING) y el Sistema Interconectado Central (SIC) en 2017 consolidó una red única que, al cierre de 2024, alcanzó una capacidad instalada de 36.778 MW \cite{CEN2024}. Dentro de este crecimiento, las Energías Renovables No Convencionales (ERNC) alcanzaron por primera vez el 50,3\% (18.522MW) de la capacidad instalada, lideradas por la energía solar con un 23\% y la eólica con un 14\%. En términos de generación, las ERNC aportaron el 40,4\% de los 85.519 GWh producidos durante 2024, lo que representa un aumento del 10,9\% respecto al año anterior \cite{CEN2024}. Esta expansión sitúa a Chile como uno de los sistemas eléctricos con mayor penetración renovable de América Latina, con una trayectoria de crecimiento que en cinco años (2019 - 2024) llevó la participación ERNC del 20\% al 40\% de la generación total \cite{LaTercera}.

Sin embargo, esta transición energética ha introducido una volatilidad estructural en la operación del sistema que no tiene precedentes en la historia del sector energético chileno. La intermitencia inherente a las fuentes solares y eólicas, combinada con una capacidad instalada solar que ya supera los 11.000 MW, suficiente para cubrir prácticamente la totalidad de la demanda del SEN, que oscila entre 10.700 y 12.200 MW, genera episodios recurrentes de sobreoferta \cite{LaTerceraPerdidas}. La consecuencia directa es la caída del costo marginal a valores cercanos a cero durante el bloque solar en zonas de alta concentración renovable, seguida de alzas abruptas en horarios de menor generación o en períodos de déficit hídrico \cite{REVE}. Este fenómeno se manifiesta concretamente en los vertimientos o curtailments. En 2024 se perdieron 5.908,71 GWh de energía eólica y solar, un aumento del 149\% respecto a los 2.375,86 GWh de 2023, equivalentes al 17\% de toda la generación renovable disponible \cite{LaTerceraPerdidas}. En el mes más crítico del año, noviembre, las plantas solares operaron menos de la mitad de su tiempo potencial por exceso de oferta.

La volatilidad resultante del costo marginal, principal señal de precio en el mercado spot del SEN, tiene implicancias directas sobre la viabilidad financiera de los agentes del mercado, la programación del despacho eléctrico y el diseño de contratos de suministro de mediano y largo plazo. A diferencia de mercados con matrices estables, el SEN combina tres fuentes de incertidumbre simultáneas: la variabilidad climática que afecta tanto a la generación solar y eólica como a los embalses hidroeléctricos, los cambios estructurales derivados de la incorporación continua de nueva capacidad renovable, con tasas de crecimiento superiores al 22\% anual frente a un crecimiento de la demanda de apenas el 2\%, y las congestiones en la red de transmisión que desacoplan los precios entre zonas del sistema \cite{Humphreys}.

El SEN es operado por el Coordinador Eléctrico Nacional (CEN) bajo un esquema liberalizado y descentralizado, que distingue segmentos regulados (transmisión y distribución) de segmentos no regulados (generación y comercialización), y abastece a aproximadamente el 98,5\% de la población chilena a lo largo de más de 3.100 km de extensión geográfica, sin interconexión con sistemas eléctricos de otros países \cite{LeonSMPChile}. Esta topología radial y extendida, combinada con una expansión de la red de transmisión que no ha logrado seguir el ritmo de la incorporación de generación renovable, ha convertido la congestión de transmisión en un fenómeno estructural y persistente antes que estocástico, a diferencia de lo observado en sistemas eléctricos densamente interconectados \cite{LeonSMPChile}.

León et al. (2026) \cite{LeonSMPChile} cuantifican esta congestión para el año 2024 mediante el porcentaje de horas en que el costo marginal de un corredor de transmisión difiere en 7\% o más del de sus nodos vecinos (\textit{decoupling}). La Tabla \ref{tab:congestion_corredores} reproduce sus resultados para los corredores de transmisión que conectan, de norte a sur, cinco de las ocho barras estudiadas en este trabajo, usando la notación original del paper, SMP (\textit{System Marginal Price}), equivalente al costo marginal (CMg) empleado en el resto de este trabajo.

\begin{table}[htbp]
    \centering
    \small
    \resizebox{\textwidth}{!}{%
    \begin{tabular}{lrrrr}
        \toprule
        & \multicolumn{2}{c}{Enero 2024} & \multicolumn{2}{c}{Julio 2024} \\
        \cmidrule(lr){2-3} \cmidrule(lr){4-5}
        Corredor & \% horas desacopladas & $\Delta$SMP (USD/MWh) & \% horas desacopladas & $\Delta$SMP (USD/MWh) \\
        \midrule
        Crucero--Cardones          & 4,0\%  & 5,1  & 5,0\%  & 5,8  \\
        Cardones--Pan de Azúcar    & 0,1\%  & 22,1 & 2,7\%  & 4,5  \\
        Pan de Azúcar--Quillota    & 1,7\%  & 10,3 & 7,8\%  & 7,6  \\
        Quillota--Alto Jahuel      & 7,9\%  & 16,3 & 5,0\%  & 46,4 \\
        Alto Jahuel--Charrúa       & 0,3\%  & 4,2  & 28,5\% & 6,9  \\
        Charrúa--Puerto Montt      & 34,3\% & 59,4 & 43,3\% & 31,9 \\
        \bottomrule
    \end{tabular}%
    }
    \caption{Porcentaje de horas con desacople de costo marginal ($\geq$7\%) y diferencial promedio de precio ($\Delta$SMP) entre barras vecinas, para los corredores de transmisión del SEN durante enero y julio de 2024. Fuente: León et al. (2026) \cite{LeonSMPChile}, a partir de datos del Coordinador Eléctrico Nacional. Alto Jahuel es un nodo intermedio entre Quillota y Charrúa que no forma parte de las ocho barras estudiadas en este trabajo.}
    \label{tab:congestion_corredores}
\end{table}

El patrón es marcadamente asimétrico: el corredor Charrúa--Puerto Montt concentra, por un margen amplio, la mayor tasa de desacople del sistema (34,3\%--43,3\% de las horas) y uno de los diferenciales de precio más altos (hasta 59,4 USD/MWh), muy por encima del resto de los corredores, que en su mayoría permanecen acoplados más del 90\% del tiempo. Esta evidencia, externa e independiente al presente trabajo, es consistente con la decisión metodológica de este trabajo de tratar a Puerto Montt de forma diferenciada del resto de las barras del sistema: la congestión estructural de su corredor de conexión, sumada a su dependencia de la disponibilidad hidroeléctrica del sur del país, sugiere que los mecanismos que determinan su costo marginal difieren cualitativamente de los que operan en las barras del norte solar, mejor conectadas entre sí.

Esta combinación hace del forecasting del precio spot en el SEN un problema de predicción de series temporales especialmente complejo, que demanda enfoques metodológicos capaces de capturar tanto la estacionalidad múltiple como las no-linealidades y los quiebres estructurales propios de un mercado en plena transición, esto motiva una revisión detallada de los modelos disponibles en la literatura.

\section{Modelos actuales}

La predicción del precio en mercados eléctricos ha sido abordada históricamente mediante modelos estadísticos clásicos como ARIMA y sus variantes estacionales, los cuales asumen estacionariedad y linealidad en las series temporales \cite{weron}. Si bien estos modelos ofrecen interpretabilidad y eficiencia computacional, su capacidad predictiva se ve limitada en mercados con alta penetración renovable como el SEN, donde la serie de costos marginales exhibe heterocedasticidad, quiebres estructurales asociados a la incorporación continua de nueva capacidad instalada, y patrones de estacionalidad múltiple determinados simultáneamente por ciclos horarios, diarios y estacionales \cite{weron}. Adicionalmente, estos modelos no capturan el comportamiento no lineal de la serie durante episodios de sobreoferta renovable, donde el costo marginal puede colapsar a cero o exhibir alzas abruptas en cuestión de horas, comportamiento que los modelos lineales son estructuralmente incapaces de representar.

Más allá de la elección de algoritmo, Lago et al. (2021) \cite{Lago2021Review} revisan el estado del arte en forecasting de precios eléctricos day-ahead y documentan que buena parte de la literatura reporta mejoras de desempeño sin verificar su significancia estadística ni evitar la fuga de información entre los conjuntos de validación y prueba usados para seleccionar el modelo final, dos prácticas que los autores identifican como necesarias para que las comparaciones entre modelos sean creíbles. Este trabajo sigue explícitamente ambas recomendaciones: la selección del mejor modelo y método de ensamblaje se realiza únicamente sobre el conjunto de validación, reportando el conjunto de prueba una sola vez (Sección 3.5), y las comparaciones entre estrategias se acompañan de un test de significancia estadística que considera la autocorrelación propia de errores horarios consecutivos (Diebold-Mariano con varianza Newey-West, Sección 4.4.4), en vez de limitarse a la diferencia nominal de MAE entre modelos.

En respuesta a estas limitaciones, la literatura ha explorado modelos de aprendizaje profundo con mayor expresividad. En particular, el \textit{Temporal Fusion Transformer} (TFT), propuesto por Lim et al. (2021) \cite{TFT}, representa un avance significativo al combinar capas recurrentes para el procesamiento local de dependencias de corto plazo con mecanismos de auto-atención interpretables para dependencias de largo alcance \cite{TFT}. A diferencia de los Transformers comunes \cite{Vaswani2017Attention}, el TFT incorpora encoders de covariables estáticas, selección de variables mediante redes de selección con compuertas, y salidas probabilísticas mediante regresión cuantílica, lo que lo hace especialmente adecuado para contextos de forecasting multihorizonte con múltiples variables exógenas, como las que caracterizan al mercado eléctrico chileno. Sin embargo, el TFT opera típicamente sobre la serie cruda sin descomponer explícitamente tendencia y estacionalidad, lo que puede dificultar su adaptación en presencia de cambios estructurales graduales como los que experimenta el SEN durante su acelerada transición energética.

Otra arquitectura Transformer relevante para forecasting de series temporales es Informer, propuesto por Zhou et al. (2021) \cite{informer}, que introduce atención \textit{ProbSparse} para reducir la complejidad cuadrática de la auto-atención estándar y así abordar horizontes de predicción y ventanas de contexto del orden de miles de pasos. Este trabajo no adopta Informer porque su aporte principal, la eficiencia computacional en secuencias muy largas, no es la restricción relevante aquí: la ventana de contexto usada por el TFT en este trabajo es de 168 pasos (una semana horaria, Sección 3.4), muy por debajo del régimen de miles de pasos que motiva a Informer. El TFT, en cambio, ofrece mecanismos --selección de variables interpretable, encoders de covariables estáticas, salidas cuantílicas nativas-- directamente alineados con la necesidad de este trabajo de incorporar múltiples variables exógenas por barra y de sustentar el análisis de interpretabilidad de la Sección \ref{sec:interpretabilidad_sc}, que Informer no provee de la misma forma.

Un aspecto técnico relevante en el diseño de arquitecturas Transformer para series temporales es el rol de las capas de normalización. La normalización de capas, \textit{Layer Normalization} (LayerNorm), ha sido considerada un componente esencial en prácticamente todas las arquitecturas Transformer desde su introducción \cite{Vaswani2017Attention}. Sin embargo, Zhu et al. (2025) \cite{Zhu2025Transformers} demuestran que es posible prescindir de estas capas mediante la introducción de \textit{Dynamic Tanh} (DyT), cuya operación central se define como $\text{DyT}(x) = \tanh(\alpha x)
$, donde $\alpha$ es un parámetro escalar aprendible que regula el escalado de la entrada antes de aplicar la función de activación. En la implementación completa, esta operación se complementa con una transformación afín posterior $\gamma \cdot \text{DyT}(x) + \beta$, donde $\beta$ y $\gamma$ son parámetros vectoriales aprendibles equivalentes a los que LayerNorm ya incorpora en su etapa de reescalado, no son parte de la operación DyT en sí, sino el mecanismo estándar que permite a la capa restaurar la escala original de las activaciones. La propuesta se inspira en la observación empírica de que las capas LayerNorm en Transformers entrenados producen mapeos entrada-salida con forma de $S$, cuyo comportamiento puede ser replicado eficientemente por la función tanh escalada. Los autores validan que los Transformers con DyT igualan o superan el rendimiento de sus contrapartes normalizadas en tareas diversas, sin requerir ajuste de hiperparámetros. La comparación entre TFT con LayerNorm y TFT con DyT constituye, por tanto, una interrogante abierta y metodológicamente relevante en el contexto de forecasting de series temporales energéticas, dado que las propiedades de estabilización del gradiente y la velocidad de convergencia pueden diferir significativamente en series con alta variabilidad y episodios de spike como las del SEN.

Los modelos híbridos representan una alternativa prometedora al combinar las fortalezas complementarias de enfoques distintos. Prophet, desarrollado por Taylor y Letham (2018) \cite{Prophet} ofrece una descomposición explícita de tendencia, estacionalidad múltiple y efectos de calendario mediante un modelo aditivo, resultando especialmente adecuado para series con patrones estacionales fuertes, valores faltantes y cambios de tendencia abruptos \cite{Prophet}. La combinación de Prophet con el TFT permitiría, en principio, aprovechar la estructura interpretable de la descomposición temporal para capturar la componente sistemática de la serie, mientras el TFT modela las dependencias no lineales. Huang et al. (2024) \cite{ProphetTransformers} exploran directamente esta combinación, acoplando Prophet con un modelo Transformer para la predicción de precios en sistemas eléctricos con alta penetración renovable, y reportan mejoras de desempeño respecto a sus modelos base considerados por separado; a diferencia del presente trabajo, no comparan esta combinación contra una estrategia alternativa sin Prophet, no evalúan variantes de normalización dentro del Transformer, ni incorporan cuantificación de incertidumbre sobre sus predicciones. No obstante, la literatura sobre modelos híbridos para mercados eléctricos latinoamericanos en particular es escasa, y los estudios existentes raramente incorporan mecanismos formales de ensamblaje, como \textit{Stacking Optimization}, ni abordan la cuantificación de incertidumbre con garantías estadísticas. Trabajos recientes ilustran esta limitación: Zhao et al. (2025) \cite{Zhao2025LIME} comparan ocho modelos de \textit{machine learning} clásico (regresión lineal, Ridge, árbol de decisión, KNN, Random Forest, SVR, XGBoost, entre otros) para el mercado eléctrico español, obteniendo buen desempeño puntual (KNN, el mejor modelo, con $R^2=0.865$) y aplicando LIME para interpretabilidad \textit{post-hoc} de las variables más influyentes, pero sin incorporar ningún mecanismo de cuantificación de incertidumbre sobre esas predicciones. Las predicciones puntuales, predominantes en la literatura de forecasting eléctrico, son insuficientes para la toma de decisiones bajo riesgo en mercados con la volatilidad estructural observada en el SEN, lo que motiva la incorporación de \textit{Conformal Prediction} (CP) como mecanismo de cuantificación de incertidumbre con cobertura garantizada sin supuestos distribucionales \cite{ConformalPrediction}.

\section{Propuesta}

En este contexto, la presente tesis propone y compara empíricamente dos estrategias de construcción de \textit{features} para un modelo híbrido de aprendizaje automático orientado a la predicción del costo marginal en el mercado eléctrico chileno. La primera, denominada Prophet+Clima (P+C), utiliza Prophet para la descomposición explícita de tendencia y estacionalidad múltiple, cuyas componentes alimentan al TFT como covariables conocidas, y ensambla las predicciones de Prophet y del TFT mediante \textit{Stacking Optimization}, un esquema de combinación que aprende los pesos óptimos de cada componente a partir de sus errores de predicción en un conjunto de validación. La segunda, denominada Sol+Clima (S+C), reemplaza las componentes de Prophet por variables de posición solar (elevación solar y horario de verano) calculadas directamente a partir de la geolocalización de cada barra, prescindiendo tanto de Prophet como del ensamblaje por Stacking. Ambas estrategias se evalúan sobre dos variantes del TFT, una con LayerNorm y otra con DyT, para la captura de dependencias temporales no lineales de largo alcance. Uno de los hallazgos centrales de este trabajo es que S+C, pese a ser la estrategia más simple, supera de forma consistente a P+C en la mayoría de las barras del sistema, lo que cuestiona la necesidad de un módulo de descomposición temporal explícito cuando se dispone de una señal física directa del ciclo diario y anual del costo marginal.

Puerto Montt, la barra del sistema con mayor nivel y dispersión histórica de costo marginal, motivó una extensión adicional: dado que su dinámica está dominada por la disponibilidad hidroeléctrica del sur del país, se evalúa si reemplazar o complementar las covariables de S+C por variables hídricas propias de la barra (cota de embalse y agua caída horaria) mejora su desempeño respecto de las estrategias generales P+C y S+C.

La incertidumbre de las predicciones se cuantifica mediante \textit{Conformal Prediction} (CP), un marco de inferencia estadística que garantiza intervalos de cobertura válidos con probabilidad preestablecida y sin imponer supuestos sobre la distribución de los errores \cite{ConformalPrediction}. Además de las variantes clásicas de CP, se propone \textit{LW-ACP} (\textit{Locally-Weighted Adaptive Conformal Prediction}), una variante que combina la ponderación local de la incertidumbre por hora del día con la adaptación \textit{online} del nivel de cobertura, aprovechando que el error del modelo en el SEN no es homogéneo a lo largo del día sino que se concentra en las horas de transición del bloque solar.

Finalmente, dado que todos los experimentos anteriores se realizan a resolución horaria, se explora de manera preliminar si la resolución temporal de los datos afecta la capacidad predictiva del modelo, comparando el desempeño de S+C a nivel diario, horario y de 15 minutos sobre un subconjunto de barras.

El modelo es evaluado sobre datos históricos del SEN, con el objetivo de determinar qué estrategia de construcción de \textit{features}, arquitectura de normalización y esquema de cuantificación de incertidumbre ofrece el mejor equilibrio entre desempeño predictivo, validez estadística y eficiencia de los intervalos resultantes. Las principales contribuciones de este trabajo son:
\begin{itemize}
    \item[(i)] La comparación empírica entre una estrategia híbrida basada en Prophet con ensamblaje formal (P+C) y una estrategia más simple basada en variables de posición solar (S+C) para el mercado eléctrico chileno, encontrándose que esta última es consistentemente superior en la mayoría de las barras del sistema,

    \item[(ii)] la comparación empírica de LayerNorm y Dynamic Tanh en un contexto de \textit{forecasting} energético,

    \item[(iii)] la incorporación de Conformal Prediction como mecanismo de incertidumbre con garantías de cobertura en el SEN, incluyendo LW-ACP, una variante que combina ponderación local por hora del día con adaptación \textit{online} del nivel de cobertura,

    \item[(iv)] un análisis diferenciado de Puerto Montt, la barra de mayor volatilidad del sistema, evaluando si variables hídricas específicas de la zona mejoran su predicción respecto de las estrategias generales, y

    \item[(v)] una exploración preliminar del efecto de la granularidad temporal de los datos (diaria, horaria y de 15 minutos) sobre el desempeño del modelo.
\end{itemize}
 
\section{Objetivos}

\subsection{Objetivo General}

El objetivo general del proyecto es desarrollar y comparar dos estrategias híbridas de aprendizaje automático, una basada en Prophet con ensamblaje formal (P+C) y otra basada en variables de posición solar (S+C), para la predicción del costo marginal en el mercado eléctrico chileno mediante el Temporal Fusion Transformer, incorporando cuantificación de incertidumbre con garantías estadísticas de cobertura mediante Conformal Prediction.

\subsection{Objetivos Específicos}

\begin{enumerate}
    \item Caracterizar el comportamiento histórico del costo marginal en el Sistema Eléctrico Nacional, identificando sus patrones de estacionalidad, tendencia y volatilidad estructural asociados a la penetración de energías renovables no convencionales.

    \item Diseñar e implementar dos estrategias de construcción de \textit{features} para el TFT, P+C (basada en la descomposición de Prophet) y S+C (basada en variables de posición solar), como modelos de captura de dependencias no lineales.

    \item Comparar empíricamente el TFT con LayerNorm y con Dynamic Tanh como mecanismo de normalización, evaluando mediante pruebas estadísticas pareadas si las diferencias de desempeño observadas entre ambas arquitecturas son sistemáticas o atribuibles al azar muestral.

    \item Desarrollar un esquema de ensamblaje mediante Stacking Optimization que combine Prophet y el TFT dentro de la estrategia P+C, aprendiendo los pesos de combinación a partir del desempeño en un conjunto de validación.

    \item Comparar estadísticamente el desempeño de P+C y S+C mediante pruebas no paramétricas pareadas, determinando si las diferencias observadas son sistemáticas o atribuibles al azar muestral.

    \item Implementar Conformal Prediction como mecanismo de cuantificación de incertidumbre, incluyendo variantes clásicas y LW-ACP, una propuesta de este trabajo que pondera la incertidumbre localmente por hora del día y adapta el nivel de cobertura en línea, y validar empíricamente la cobertura y eficiencia de los intervalos de predicción resultantes.

    \item Evaluar si la incorporación de variables hídricas propias de Puerto Montt mejora su desempeño respecto de las estrategias generales P+C y S+C, dada la mayor volatilidad estructural de esa barra.

    \item Explorar de manera preliminar el efecto de la granularidad temporal de los datos, comparando el desempeño de la mejor estrategia a nivel diario, horario y de 15 minutos.
\end{enumerate}

\section{Organización del documento}

El presente documento se organiza en cuatro capítulos adicionales a esta Introducción. El Capítulo 2 presenta el marco teórico: los fundamentos de series temporales, Prophet, el Temporal Fusion Transformer y la comparación LayerNorm/Dynamic Tanh, Stacking Optimization, y Conformal Prediction, incluyendo las variantes clásicas y LW-ACP. El Capítulo 3 detalla la metodología: la fuente y el preprocesamiento de los datos, la construcción de las estrategias P+C y S+C, la configuración de Prophet y del TFT, el esquema de Stacking Optimization, la variante DyT, y la implementación de las distintas variantes de Conformal Prediction. El Capítulo 4 presenta los resultados: el desempeño comparado de P+C y S+C sobre las ocho barras del SEN, el análisis de errores extremos, los resultados de Conformal Prediction, el análisis de interpretabilidad del TFT, el estudio diferenciado de Puerto Montt, y la exploración preliminar de granularidad temporal. Finalmente, se presentan las conclusiones del trabajo, sintetizando los hallazgos principales y proponiendo líneas de trabajo futuro.